\section{Conclusion} \label{sec:conclusion}

\subsection{Summary of the Project}

This project addressed a recurring practical problem in university computer labs: instructors and lab assistants cannot reliably see, in real time, which students are struggling and which questions are causing the most difficulty across a session of several dozen to over a hundred concurrent learners.
The resolution is a React and FastAPI dashboard, V2 in the staged Streamlit-then-React evolution recorded in Chapter~\ref{sec:implementation}, that ingests live submission data, computes interpretable struggle and difficulty scores from that data, and surfaces both as ranked leaderboards alongside a mistake-cluster view, a collaborative-filtering recommender, a saved-session archive, Item Response Theory and Bayesian Knowledge Tracing layers for advanced modelling, and a Retrieval-Augmented Generation feedback panel.
A second route on the same FastAPI process hosts the mobile lab-assistant interface, so coordination between the instructor and the floor team runs through a single shared session state.

The seven objectives stated in Section~\ref{sec: intro} were addressed in order.
The literature review in Chapter~\ref{sec:requirements} identified the dashboard space and surfaced the gaps in real-time, lab-scale, instructor-facing analytics; the same chapter formalised the functional and non-functional requirements that drove the rest of the project.
Chapter~\ref{sec:design} set out the data model that turns raw submission attempts into the per-student and per-question signals the scorers consume, and it defined the system architecture - shared-state coordination, view routing, caching - that Chapter~\ref{sec:implementation} then implemented.
The struggle and difficulty composites were specified in Chapter~\ref{sec:design} and refit against second-opinion labels from a large language model rater in Section~\ref{sec:eval-results}, covering composite weights, scalar hyperparameters, and band thresholds.
The dashboard itself was built incrementally from the V1 Streamlit reference to the V2 React and FastAPI implementation documented in Chapter~\ref{sec:implementation}.
Functional and non-functional evaluation followed in Chapter~\ref{sec:evaluation}, combining a retrospective replay of the deployment cohort with an eleven-question instructor and student survey.

Against the rater's four-band rating, the deployed composite achieved Spearman $\rho = +0.588$ on the struggle target and $\rho = +0.468$ on difficulty (Section~\ref{sec:eval-results}); after a constrained cross-validated brute-force search the deployed cutpoints lifted linear-weighted Cohen's $\kappa$ to $+0.384$ for struggle and $+0.387$ for difficulty (Section~\ref{sec:eval-thresholds}).
The 2PL Item Response Theory difficulty model was retained as a complementary toggle but not promoted to be the deployed scorer: its band agreement against the rater was marginally negative ($\kappa = -0.024$ at trained thresholds) and the top-ten hardest set overlapped with the composite's by one question in ten (Section~\ref{sec:eval-irt}).
The system is deployed and in use; the sections that follow set out what the validation contributed and what remains for future work.

\subsection{Key Findings and Contributions}

The first contribution is a systematic empirical validation of every tunable component of the deployed scoring pipeline against second-opinion labels from a large language model rater, with positive and negative findings reported in equal weight.
The validation covered three tiers of parameter.
The composite weights were refit by ordinary least-squares regression against the four-band rater target (Section~\ref{sec:eval-results}); the two scalar hyperparameters, the Bayesian shrinkage constant $K$ and the collaborative-filtering threshold $\tau$, were tuned by Tree-structured Parzen Estimator across fifty-trial Optuna studies (Section~\ref{sec:eval-hyperopt}); and the score-to-band thresholds were trained by brute-force linear-weighted $\kappa$ maximisation across six model variants under constrained group five-fold and leave-one-out cross-validation (Section~\ref{sec:eval-thresholds}).

The positive findings are concrete.
Tuning $\tau$ from the hand-set $0.7$ to the tuned $0.900$ raised Spearman $\rho$ by $+0.160$, a substantial gain indicating the initial setting was too permissive.
Four of the six threshold variants cleared the pre-registered promotion bar of $\Delta\kappa \geq +0.05$ under constrained cross-validation, and a follow-up search on the deployed min-max composite produced the cutpoints now in \texttt{config.py} with band-agreement $\kappa = +0.384$ for struggle and $\kappa = +0.387$ for difficulty.
The negative findings are reported with the same emphasis.
The shrinkage constant $K$ moved from the hand-set $5$ to the tuned $0$ but the gain $\Delta\rho = +0.009$ sat within per-fold standard deviation, validating the hand-set value rather than refining it.
The 2PL IRT difficulty model failed to clear the band-agreement bar at any threshold pass, with $\kappa = -0.016$ at hand-set cutpoints and $\kappa = -0.024$ at trained cutpoints, and was retained as a complementary view rather than promoted to the deployed scorer (Section~\ref{sec:eval-irt}).
The contribution is the systematic coverage and the documented negative findings.

The second contribution is the training methodology itself, reusable beyond this dashboard.
The methodology has three components.
The supervision target is a four-band ordinal rating (On Track / Minor Issues / Struggling / Needs Help on struggle; Easy / Medium / Hard / Very Hard on difficulty) produced by a large language model rater shown per-snapshot and per-question signals.
The estimator is ordinary least-squares over the composite features rather than gradient boosting or random forest.
The evaluation metric stack is Spearman $\rho$ for ranking, linear-weighted Cohen's $\kappa$ for band agreement, and mean absolute error for calibration, replacing the binary-only AUC of the original framing.

The model-class-selection rigour is the bake-off in Section~\ref{sec:eval-results}: ordinary least-squares regression was benchmarked against random forest and gradient boosting across all three targets under matched cross-validation.
Ordinary least-squares won the struggle target decisively at $\rho = +0.588$ against the next-best random forest at $\rho = +0.565$, tied the difficulty target within per-fold noise at $\rho = +0.468$ against random forest's $\rho = +0.479$, and was the only competitive class with interpretable signed weights.
Tree based alternatives offered no consistent gain large enough to outweigh that interpretability and were dropped from the deployed path.


\subsection{Future Work}

\paragraph{Cross-session aggregation}
The deployed dashboard reports struggle and difficulty within a single live session.
A respondent to the closing survey question (Section~\ref{sec:eval-survey}, Q11) asked for a view of how frequently a given student is flagged as struggling across multiple lab sessions, not just within one.
Implementing this requires persisting per-snapshot scoring beyond the live session window and aggregating per-student across sessions in a way that respects the session-grouped cross-validation structure used in evaluation.
This is the most direct survey-driven feature request in the chapter.

\paragraph{Joint weight and threshold optimisation via ordinal logistic regression}
The pipeline currently trains composite weights (Section~\ref{sec:eval-results}) and band thresholds (Section~\ref{sec:eval-thresholds}) in two separate passes; the recalibration pass on the deployed min-max composite was needed because the originally-promoted cutpoints were fit against a different score scale than the one the deployed system serves.
A proportional-odds ordinal logistic regression, fit either through \texttt{mord} or \texttt{statsmodels.OrderedModel}, would optimise the weights and the cutpoints jointly against the four-band ordinal target in a single fit, removing the scale-mismatch failure mode at its root rather than recalibrating around it.

\paragraph{Cohort-balanced re-evaluation}
The validation cohort produced no questions in the rater's "Easy" band and only $5.8\%$ of snapshots in the "On Track" band; the cohort distributions reported in Section~\ref{sec:eval-results} reflect a cohort skewed toward higher difficulty and lower performance which may be due to no intervention.
The lowest band of each scorer is therefore effectively empty on this cohort, so band-agreement findings against the rater are conditioned on a distribution the next cohort may not replicate.
Re-running the evaluation on a more balanced cohort, or on stratified sub-cohorts of the same module, would test whether the recalibrated cutpoints generalise and whether the IRT difficulty anti-result holds on a less-skewed response matrix.

\paragraph{Per-semester threshold re-tuning as a maintenance concern}
Following directly from the cohort caveat in Section~\ref{sec:eval-thresholds}, the threshold-training pass should be re-run per module rather than the existing tuples reused indefinitely.
This is a deployment-side machine-learning maintenance concern rather than a model-design change: the underlying composite score and threshold is not portable across cohorts, and the maintenance loop should recognise that.

\paragraph{Restoring the OpenAI incorrectness scorer to a non-fallback rate}
Section~\ref{sec:eval-tradeoffs} flags this as the highest-leverage future change in the evaluation chapter, ranking above retraining any of the weights or thresholds.
The OpenAI incorrectness scorer fell back to a midpoint score on $92.5\%$ of the production submissions, which propagates into the composite as noise on the strongest single signal in the weight vector.
A more robust scorer path: retry logic with backoff, structured-output validation, or a local fallback model rather than a midpoint imputation would raise the ceiling on every downstream model in the pipeline.

\paragraph{Prospective intervention study}
The retrospective replay design used in Chapter~\ref{sec:evaluation} explicitly scopes out a prospective intervention study, in which dashboard-informed instructor decisions are compared with a control group on student outcomes.
A future evaluation pass would need a controlled comparison, a lift of the student-identifier anonymisation that currently rules out grade linkage, and a submission-to-lab-seat mapping that does not exist in the current data path.
This is named here as the next-generation evaluation step rather than as an immediate follow-up; it is the only way to attribute outcome changes to the dashboard rather than to correlations with it.

\paragraph{Smart device integration}
The original scope in Section~\ref{sec: intro} named smartwatch and smart-glasses notifications to lab assistants as a secondary aim.
The V2 deployment scoped this down to a mobile lab-assistant web route served by the same FastAPI process, which delivers the same coordination affordance through a web interface rather than a native device.
Future work would restore the original goal with native notifications to a wearable, layered over the same file-locked shared-state coordination V2 already uses for cross-stack session synchronisation.
This closes the loop on the secondary aim from Chapter~\ref{sec: intro} rather than the modelling-side findings that drive the earlier items.